\chapter{Vitrectomy}
\section*{Overview}
Vitrectomy removes some or all of the vitreous gel from inside the eye and replaces it with fluid, gas, or oil as needed.
\section*{Common Instruments and Equipment}
The exact setup varies by technique and setting. Common items may include:
\begin{itemize}[nosep,leftmargin=*]
  \item operating microscope and small-gauge vitrectomy system
  \item intraocular illumination and infusion equipment
  \item laser probe and gas or silicone oil when required
\end{itemize}

\section*{Why It Is Done}
It may treat retinal detachment, vitreous hemorrhage, macular hole, epiretinal membrane, endophthalmitis, or complications of eye trauma.
\section*{How It Is Performed}
Through tiny ports, microsurgical instruments remove vitreous and address the retinal or macular problem. Laser, cryotherapy, gas, or silicone oil may be used.
\section*{Contraindications}
No single contraindication applies to every patient; the treating team assesses whether the procedure is suitable for the individual and the reason for treatment.

\section*{Precautions}
The care team reviews relevant health conditions, medicines, allergies, and procedure-specific risks, and explains any changes needed for the individual.

\section*{Risks and Recovery}
Infection, bleeding, increased eye pressure, cataract progression, retinal redetachment, and vision loss are possible. Positioning and restrictions are critical when a gas bubble is used; air travel may be prohibited.
\section*{References}
\begin{enumerate}\item MedlinePlus. Vitrectomy. U.S. National Library of Medicine; 2025.\item Current Medical Diagnosis and Treatment. McGraw-Hill; 2025.\end{enumerate}
\clearpage
